\section{Введение}

Лабораторная работа посвящена исследованию режимов возбуждения лампового генератора и зависимости амплитуды и частоты автоколебаний от параметров системы. Теоретическое описание выполняется с использованием метода медленно меняющихся амплитуд (метода Ван-дер-Поля), а полученные результаты сопоставляются с исследованием фазовой плоскости генератора.

Исследуемая система представляет собой квазисинусоидальную автоколебательную систему. В зависимости от положения рабочей точки на характеристике лампы возможны мягкий, жёсткий и сложно-жёсткий режимы возбуждения.

\subsection*{Цель} \addcontentsline{toc}{subsection}{Цель}

Целью работы является изучение различных режимов возбуждения лампового генератора и установление зависимости амплитуды и частоты автоколебаний от параметров системы.

\subsection*{Задания} \addcontentsline{toc}{subsection}{Задания}

\begin{enumerate}
	\item Познакомиться с методом медленно меняющихся амплитуд.
	\item Составить исходные и укороченные уравнения исследуемой системы.
	\item Найти состояния равновесия исходных и укороченных уравнений, исследовать их устойчивость и определить их типы.
	\item Теоретически построить бифуркационные диаграммы и фазовые портреты для различных режимов работы.
	\item Познакомиться с экспериментальной установкой.
	\item Экспериментально снять бифуркационные диаграммы зависимости амплитуды автоколебаний от коэффициента взаимоиндукции при мягком и сложно-жёстком режимах возбуждения.
	\item Сопоставить экспериментальные результаты с результатами теоретического анализа.
\end{enumerate}

\subsection*{Оборудование} \addcontentsline{toc}{subsection}{Оборудование}

\begin{itemize}
	\item Лабораторная установка лампового генератора (с регулируемыми напряжением смещения сетки и коэффициентом взаимоиндукции).
	\item Электронно-лучевой осциллограф (для наблюдения фазовой плоскости).
	\item Вольтметры.
	\item Источники питания (анодное напряжение и напряжение смещения сетки).
	\item Соединительные провода.
\end{itemize}
